% GENERATED FILE. Edit canonical lessons, then run npm run generate:books.
% canonical-source-hash: fnv1a64:6a1fa75f605cc362
% canonical-lessons: BN-C59-ruti, BN-C59-am, BN-C59-sobji, BN-C59-bari, BN-C59-janala

\chapter{Bread, Mango, Vegetables, House, Window}
\label{ch:bn-bread-mango-vegetables-house-window}
\hlchaptermodality{59}

\begin{chapteropening}
\textbf{By the end of this chapter:} \emph{I can say bread, a mango, vegetables, a house, home and a window.}

Complete the last lesson of chapter 59: I can say bread, a mango, vegetables, a house, home and a window.
\end{chapteropening}

\section[\texorpdfstring{\bn{রুটি} (ruṭi)}{রুটি (ruṭi)}]{\bn{রুটি} (ruṭi) — bread}

\label{lesson:BN-C59-ruti}



\begin{quote}
Before the new one: say the Bengali for yoghurt, then the Bengali for an onion.
\end{quote}

\subsection*{You'll want to know: \bn{রুটি}}
\textbf{\bn{রুটি}} — \emph{ruṭi} — \textquotedblleft{}bread\textquotedblright{}.

\begin{cousinweb}
\textbf{\bn{রুটি}} (\emph{ruṭi}) is a tadbhava, a word inherited and reshaped by sound change. It comes from \emph{*roṭṭa}, \textquotedblleft{}bread\textquotedblright{}, a form CDIAL reconstructs from its descendants (CDIAL 10837). Late Sanskrit has \emph{rōṭikā}.
\end{cousinweb}

\begin{grammarlens}[title={What you've built}]
One more word for this chapter.
\end{grammarlens}

\subsection*{Your turn}
\begin{itemize}
\raggedright
  \item \emph{Say it:} \emph{ruṭi}
  \item \emph{Say it:} \emph{ruṭi}, once more
  \item \emph{Say it:} say \emph{ruṭi}
  \item \emph{Recall:} say the Bengali for yoghurt, then the Bengali for an onion, then say \emph{ruṭi} again
\end{itemize}

\begin{tcolorbox}[breakable,colback=teal!4,colframe=teal!35!black,title={Before you move on}]
What does \bn{রুটি} mean? (\textquotedblleft{}Bread\textquotedblright{}.) Say it once more.
\end{tcolorbox}
\section[\texorpdfstring{\bn{আম} (ām)}{আম (ām)}]{\bn{আম} (ām) — a mango}

\label{lesson:BN-C59-am}



\begin{quote}
Before the new one: say the Bengali for an onion, then the Bengali for bread.
\end{quote}

\subsection*{You'll want to know: \bn{আম}}
\textbf{\bn{আম}} — \emph{ām} — \textquotedblleft{}a mango\textquotedblright{}.

\begin{cousinweb}
\textbf{\bn{আম}} (\emph{ām}) is a tadbhava, a word inherited from Sanskrit and reshaped by sound change. It comes from \emph{āmra}, \textquotedblleft{}mango tree; its fruit\textquotedblright{} (CDIAL 1268).
\end{cousinweb}

\begin{grammarlens}[title={What you've built}]
One more word for this chapter.
\end{grammarlens}

\subsection*{Your turn}
\begin{itemize}
\raggedright
  \item \emph{Say it:} \emph{ām}
  \item \emph{Say it:} \emph{ām}, once more
  \item \emph{Say it:} say \emph{ām}
  \item \emph{Recall:} say the Bengali for an onion, then the Bengali for bread, then say \emph{ām} again
\end{itemize}

\begin{tcolorbox}[breakable,colback=teal!4,colframe=teal!35!black,title={Before you move on}]
What does \bn{আম} mean? (\textquotedblleft{}A mango\textquotedblright{}.) Say it once more.
\end{tcolorbox}
\section[\texorpdfstring{\bn{সবজি} (sôbji)}{সবজি (sôbji)}]{\bn{সবজি} (sôbji) — vegetables}

\label{lesson:BN-C59-sobji}



\begin{quote}
Before the new one: say the Bengali for bread, then the Bengali for a mango.
\end{quote}

\subsection*{You'll want to know: \bn{সবজি}}
\textbf{\bn{সবজি}} — \emph{sôbji} — \textquotedblleft{}vegetables\textquotedblright{}.

\begin{cousinweb}
\textbf{\bn{সবজি}} (\emph{sôbji}) is a loan, usually traced to Persian \emph{sabzī}, \textquotedblleft{}greenness, greens\textquotedblright{}. It is built on Persian \emph{sabz}, \textquotedblleft{}green\textquotedblright{}.
\end{cousinweb}

\begin{grammarlens}[title={What you've built}]
One more word for this chapter.
\end{grammarlens}

\subsection*{Your turn}
\begin{itemize}
\raggedright
  \item \emph{Say it:} \emph{sôbji}
  \item \emph{Say it:} \emph{sôbji}, once more
  \item \emph{Say it:} say \emph{sôbji}
  \item \emph{Recall:} say the Bengali for bread, then the Bengali for a mango, then say \emph{sôbji} again
\end{itemize}

\begin{tcolorbox}[breakable,colback=teal!4,colframe=teal!35!black,title={Before you move on}]
What does \bn{সবজি} mean? (\textquotedblleft{}Vegetables\textquotedblright{}.) Say it once more.
\end{tcolorbox}
\section[\texorpdfstring{\bn{বাড়ি} (bāṛi)}{বাড়ি (bāṛi)}]{\bn{বাড়ি} (bāṛi) — a house, home}

\label{lesson:BN-C59-bari}



\begin{quote}
Before the new one: say the Bengali for a mango, then the Bengali for vegetables.
\end{quote}

\subsection*{You'll want to know: \bn{বাড়ি}}
\textbf{\bn{বাড়ি}} — \emph{bāṛi} — \textquotedblleft{}a house, home\textquotedblright{}.

\begin{cousinweb}
\textbf{\bn{বাড়ি}} (\emph{bāṛi}) is a tadbhava, a word inherited from Sanskrit and reshaped by sound change. It comes from \emph{vāṭikā}, \textquotedblleft{}enclosure, garden\textquotedblright{} (CDIAL 11480). CDIAL glosses the Bengali form \textquotedblleft{}garden\textquotedblright{}.
\end{cousinweb}

\begin{grammarlens}[title={What you've built}]
One more word for this chapter.
\end{grammarlens}

\subsection*{Your turn}
\begin{itemize}
\raggedright
  \item \emph{Say it:} \emph{bāṛi}
  \item \emph{Say it:} \emph{bāṛi}, once more
  \item \emph{Say it:} say \emph{bāṛi}
  \item \emph{Recall:} say the Bengali for a mango, then the Bengali for vegetables, then say \emph{bāṛi} again
\end{itemize}

\begin{tcolorbox}[breakable,colback=teal!4,colframe=teal!35!black,title={Before you move on}]
What does \bn{বাড়ি} mean? (\textquotedblleft{}A house, home\textquotedblright{}.) Say it once more.
\end{tcolorbox}
\section[\texorpdfstring{\bn{জানালা} (jānālā)}{জানালা (jānālā)}]{\bn{জানালা} (jānālā) — a window}

\label{lesson:BN-C59-janala}



\begin{quote}
Before the new one: say the Bengali for vegetables, then the Bengali for a house.
\end{quote}

\subsection*{You'll want to know: \bn{জানালা}}
\textbf{\bn{জানালা}} — \emph{jānālā} — \textquotedblleft{}a window\textquotedblright{}.

\begin{cousinweb}
\textbf{\bn{জানালা}} (\emph{jānālā}) is a loan from Portuguese \emph{janela}, \textquotedblleft{}window\textquotedblright{}.
\end{cousinweb}

\begin{grammarlens}[title={What you've built}]
One more word for this chapter.
\end{grammarlens}

\subsection*{Your turn}
\begin{itemize}
\raggedright
  \item \emph{Say it:} \emph{jānālā}
  \item \emph{Say it:} \emph{jānālā}, once more
  \item \emph{Say it:} say \emph{jānālā}
  \item \emph{Recall:} say the Bengali for vegetables, then the Bengali for a house, then say \emph{jānālā} again
\end{itemize}

\begin{tcolorbox}[breakable,colback=teal!4,colframe=teal!35!black,title={Before you move on}]
What does \bn{জানালা} mean? (\textquotedblleft{}A window\textquotedblright{}.) Say it once more.
\end{tcolorbox}
